\section{Local markers and periodic continuations}\label{marker:section}

The clock provides algebraic control of the two tests. We now arrange
letter positions so that those tests can read enough information on
opposite sides of a cut. A local priority rule distinguishes two cases:
a nearby marker offers a boundary; the absence of nearby markers forces
a periodic block. Overlap will make many such blocks part of one common
repetition.

This is the finite-word marker principle associated with Krieger's
marker lemma~\cite[Lemma~2]{Krieger1982}. Related priority constructions
are used in \cite[Section~4]{HeightThirteen} and
\cite[Section~4.1]{HeightFour}; a general clopen merge formulation is
presented in \cite[Lemmas~3.2--3.4]{Meyerovitch2025}. We prove the exact
finite-window and overlap statements used here. They do not require an
embedding or a change of alphabet.

Positions of letters are integers. Intervals of possible marker
positions are inclusive integer intervals; a block of letters on
$[a,b)$ has length $b-a$. A positive integer $p$ is a \emph{period}
of a finite block if its letters at positions differing by $p$ agree
whenever both positions belong to the block.

\begin{lemma}[Local markers]\label{marker:local}
Fix a nonempty finite alphabet $\Sigma$ and integers $d\ge2$ and
$K>3d^2$. There are an integer $r\ge K$ and a fixed rule assigning
marker positions to every two-sided word in $\Sigma^{\mathbb Z}$
such that:
\begin{enumerate}[label=\textup{(\roman*)}]
\item distinct markers have distance at least $d$;
\item the decision at $x$ uses only the letters on $[x-r,x+r)$, and
the rule commutes with all integer translations;
\item if $[z-d,z+d]$ contains no marker, the block $[z,z+K)$ has a
positive period less than $d$.
\end{enumerate}
\end{lemma}

\begin{proof}
Order the length-$K$ words with no period below $d$ as
$v_1,\ldots,v_N$. At stage $j$, place a marker at each start of $v_j$
that is at distance at least $d$ from every earlier marker. Keep all
earlier markers. Two newly placed starts at distance $a$ with
$0<a<d$ would give two overlapping copies of $v_j$; equality on the
overlap says exactly that $a$ is a period of $v_j$. This is excluded by
the list. Thus both old--new and new--new pairs satisfy the separation
requirement at every stage.

The procedure has a finite inspection radius despite the infinite word.
Take $r_0=K$ and $r_j=K+j(d-1)$. To decide whether to place a marker at
$x$ at stage $j$, inspect its length-$K$ block and the earlier marker
status at the $2d-1$ positions from $x-(d-1)$ to $x+(d-1)$. Inductively,
the necessary windows fit inside $[x-r_j,x+r_j)$. The final radius
$r=r_N$ works; when the list is empty, $r=K$ works for the empty marker
set. Every instruction is stated in relative positions, so it commutes
with translations.

If the block beginning at $z$ appeared in the list, its start either
received a marker at its stage or was prevented by a marker at distance
less than $d$. In either event that marker remains in the final set.
Absence throughout $[z-d,z+d]$ therefore implies that the block was not
listed, which means that it has a positive period below $d$.
\end{proof}

A later search may inspect several possible origins. It is not enough
to know that each corresponding seed is periodic: they must describe the
same continuation. The following elementary overlap estimate gives this
agreement and also converts an interval without markers into a single
periodic run.

\begin{lemma}[Common periodic continuation]\label{marker:periodic}
Two two-sided words with respective positive periods $a,b$ that agree
on $ab$ consecutive letter positions agree everywhere. Consequently,
for $d\ge2$ and $K>3d^2$:
\begin{enumerate}[label=\textup{(\roman*)}]
\item a length-$K$ block having a period below $d$ determines a
unique two-sided periodic continuation with some period below $d$;
\item any finite nonempty family of such blocks in one word, whose
start positions have diameter $w$ with $K-w\ge d^2$, determines
one common continuation, agreeing with the word on their union;
\item for the rule in Lemma~\ref{marker:local}, if the inclusive
interval $[z,y]$ contains no marker and $y-z\ge2d$, the word agrees
with one two-sided word of period less than $d$ throughout
\begin{equation}\label{marker:markerless-run}
                          [z+d,y-d+K).
\end{equation}
\end{enumerate}
The same assertions apply to a finite word when all stated blocks
and marker inspection windows are present in that word.
\end{lemma}

\begin{proof}
The integer $ab$ is a period of both two-sided words. Any position
can be moved into an interval of $ab$ consecutive positions by adding
a multiple of $ab$, without changing either letter. Agreement on
that interval therefore proves agreement everywhere.

A finite block of period $a<d$ has length greater than $a$ and
extends by repeating its first $a$ letters in both directions. Any
two choices of periods $a,b<d$ yield continuations agreeing on the
block, whose length exceeds $ab$. The first assertion identifies the
continuations and proves~\textup{(i)}. For~\textup{(ii)}, any two
blocks overlap in at least $K-w\ge d^2>ab$ positions, where $a,b$
are their short periods. Their continuations are equal. Fixing any
one block identifies a common continuation for the whole family.
Since $w<K$, the union of the blocks is an interval, and all its
letters agree with this continuation.

For~\textup{(iii)}, each integer start $x$ from $z+d$ through $y-d$
has $[x-d,x+d]\subseteq[z,y]$. Lemma~\ref{marker:local} therefore
gives a period below $d$ for its length-$K$ block. Successive blocks
overlap in $K-1>d^2$ letters, so the first assertion identifies their
continuations one after another. Their union is exactly
\eqref{marker:markerless-run}. The common continuation has, for
example, the short period of the first block. This argument also
covers $y-z=2d$, when there is just one block.

Finally, extend a finite word arbitrarily to a two-sided word, which
is possible because $\Sigma$ is nonempty. Internal inspection
windows give the same marker decisions for every such extension;
the blocks and the equalities just proved use the given letters.
Thus the conclusions about the original finite word are independent
of the extension.
\end{proof}

We adopt a fixed convention for finite arguments. A required marker
inspection is performed only when its entire window lies in that
argument; otherwise the whole test rejects. In particular, an unavailable
window never removes an origin from a comparison or uniqueness test.
This convention is what permits the two-sided proof to establish facts
about the actual letters read by a prefix or suffix.
